\chapter{Artikel}
\label{cha:Artikel}

Artikel begleiten das Nomen und bestimmen seinen Kasus, Numerus und Genus.

\section{Bestimmter Artikel}
\label{sec:BestimmterArtikel}

\begin{center}
\begin{tabular}{|l|c|c|c|c|}
\hline
\textbf{Kasus} & \textbf{Maskulin} & \textbf{Feminin} & \textbf{Neutral} & \textbf{Plural} \\ \hline
Nominativ & der & die & das & die \\ \hline
Akkusativ & den & die & das & die \\ \hline
Dativ & dem & der & dem & den \\ \hline
Genitiv & des & der & des & der \\ \hline
\end{tabular}
\end{center}

\section{Unbestimmter Artikel}
\label{sec:UnbestimmterArtikel}

\begin{center}
\begin{tabular}{|l|c|c|c|c|}
\hline
\textbf{Kasus} & \textbf{Maskulin} & \textbf{Feminin} & \textbf{Neutral} & \textbf{Plural} \\ \hline
Nominativ & ein & eine & ein & keine \\ \hline
Akkusativ & einen & eine & ein & keine \\ \hline
Dativ & einem & einer & einem & keinen \\ \hline
Genitiv & eines & einer & eines & keiner \\ \hline
\end{tabular}
\end{center}

\section{Nullartikel}
\label{sec:Nullartikel}

Der Nullartikel wird verwendet bei:
\begin{itemize}
    \item \textbf{Abstrakta}: Glück, Freude, Liebe
    \item \textbf{Stoffnamen}: Wasser, Brot, Fleisch
    \item \textbf{Berufen/Nationalitäten}: Er ist Arzt. Sie ist Deutsche.
    \item \textbf{plurale Nomen ohne Artikel}: Kinder, Bücher
    \item \textbf{fixen Wendungen}: Hunger haben, Angst haben
\end{itemize}

\section{Gebrauch des Artikels}
\label{sec:ArtikelGebrauch}

\subsection{Bestimmt vs. Unbestimmt}
\begin{itemize}
    \item \textbf{Bestimmt}: Etwas Bekanntes, bereits Erwähntes, Einzigartiges
    \item \textbf{Unbestimmt}: Etwas Neues, Unbekanntes
\end{itemize}

\section{Übungen}
\label{sec:ArtikelUebungen}

\begin{enumerate}
    \item Ich brauche \underline{\hspace{1cm}} Buch. (ein)
    \item Das ist \underline{\hspace{1cm}} Haus. (das)
    \item \underline{\hspace{1cm}} Kind spielt im Garten. (ein)
    \item Ich kenne \underline{\hspace{1cm}} Mann nicht. (der)
    \item Sie hat \underline{\hspace{1cm}}Katze. (eine)
\end{enumerate}

\chapterend